\section{问题与动机}
\label{sec:problem}
\textbf{问题。} 候选解 $W=(w_1,w_2,w_3)$ 为每栋建筑给出一个权重向量 $w_b\in\Delta^6$（非负、和为一）。仿真器返回
\begin{equation}
R(W)=\frac{1}{8}\sum_{k=1}^{8}\sum_{t=2}^{168}\bigl|E_{k,t}(W)-E_{k,t-1}(W)\bigr|\ \text{kWh},
\label{eq:R}
\end{equation}
其中 $E_{k,t}$ 是第 $k$ 周第 $t$ 小时的区域净负荷。最佳训练行（\texttt{p01969}）的 $R=115.76$，这也是离线门槛；课程二次岭回归
基线的最优解实测为 119.09。

\textbf{代理模型在远离数据处失效。} 在冻结的十划分评估规程上，我们的代理模型（二次岭回归、ExtraTrees、梯度提升、堆叠集成）
得分相差都在一个标准差之内。然而它们的最优解远离数据，而我们能在那里测到的两个点都与之矛盾（表~\ref{tab:offline}）：基线
最优解预测为 75.9，实测为 119.09；我们的第一次提交是三个树模型的共识，预测为 107--111，实测为 119.71。因此第二次提交只是
从最佳行出发、长度为 0.05 的一步，由加性 ARD 高斯过程~\cite{rasmussen}在 19,483 个随机方向中选出它最有把握的方向。它实测为
115.68，比门槛低 0.08~kWh，但增益（$-0.08$）比预测小 35 倍。在线阶段同样如此：加性 GP 把 6,548 个本地实测点中最低的 200 个
（$R\approx105.1$）预测为 $124.4\pm8.8$，因为超出约两个训练点间距后，它的先验均值会回落到数据平均值。

\begin{table}[t]
\centering
\caption{第二次离线提交之前实测的两个点上的代理模型预测（kWh）。}
\label{tab:offline}
\begin{tabular}{lrr}
\toprule
模型 & 基线最优解 & 第一次提交\\
\midrule
二次岭回归（课程） & 75.90 & 95.13\\
HistGradientBoosting & 117.62 & 107.80\\
XGBoost & 114.21 & 107.28\\
堆叠集成 & 118.73 & 111.42\\
\midrule
\textbf{实测} & \textbf{119.09} & \textbf{119.71}\\
\bottomrule
\end{tabular}
\end{table}

\textbf{实测揭示的结构。} 离线阶段结束后仿真器被发布，助教允许用本地运行来设计策略，前提是提交的代码在首次调用前只使用
训练行。一次探索性搜索（约 22,500 次本地评估，从未作为策略的输入）发现了三条规律。第一，在同一个盆地内（即围绕一种主导目标
组合的低 $R$ 区域）$R$ 在建筑之间是加性的（交互误差约 0.03~kWh）。训练行在全局上也与此一致，但更宽松：按建筑加性的模型
交叉验证效果与不受约束的模型相当（$R^2$ 0.846 对 0.844），说明加性几乎不损失拟合。
第二，$R$ 呈阶梯状：沿一个权重方向，它在宽度为权重 0.02--0.2\,\% 的平台上保持不变，因此许多移动会返回完全相同的值。第三，
极小权重起开关作用：在原本不用的目标上设 $10^{-5}$ 到 $3\times10^{-3}$ 的份额，会翻转几乎持平的逐小时决策，而最后的改进来自
单坐标扫描。加性是第~\ref{sec:block}~节的动机；阶梯结构和各类移动收益的不均衡是第~\ref{sec:alloc}~节的动机。
